% Ported verbatim from paper/top3_joint_moment_bound.tex lines 1348-1384.
% Only cross-reference wording was adapted; no mathematics was altered.
\begin{proof}[Proof of Theorem~\ref{main}]
Any Borel law on $\R^3$ can be realized on a standard nonatomic probability
space, by adjoining an independent uniform variable if necessary.
For its isotropic $X$, the coordinates are an orthonormal three-frame and
$W_X=\mathbf1_{\{X\cdot X'>0\}}$ is a measurable symmetric binary kernel.
Therefore
\[
 \frac{\E|X\cdot X'|+|\E X|^2}{2}
   =\iint W_X X\cdot X'
   \le\Sigma_3(W_X)\le\mathcal M\le c.
\]
This proves~\eqref{law-bound}, including zero vectors and unbounded radii.
The radial restriction was used only for a maximizing canonical law.

For a finite graph, the exact-rank Ky Fan formula reads
\[
 \Lambda_3(G)=\max_{Q=Q^T=Q^2,\ \tr Q=3}\tr(A_GQ).
\]
Factor $Q=YY^T$ with orthonormal columns, let $y_i$ be its rows, and let
$X$ be uniform on $\sqrt n\,y_1,\ldots,\sqrt n\,y_n$. Then
\[
 \E XX^T=I_3,\qquad
 \E|X\cdot X'|=\|Q\|_1/n,\qquad
 |\E X|^2=\boldsymbol1^TQ\boldsymbol1/n,
\]
where $\|Q\|_1$ denotes the entrywise absolute sum.
Since $Q_{ii}\ge0$ and $\sum_iQ_{ii}=3$, the zero adjacency diagonal gives
\[
 \tr(A_GQ)\le\sum_{i\ne j}(Q_{ij})_+
 =\frac{\|Q\|_1+\boldsymbol1^TQ\boldsymbol1}{2}-3
 \le cn-3.
\]
Taking the maximum proves~\eqref{main-bound}. No positivity assumption on
$\lambda_3(G)$ is used; repeated eigenvalues, zero rows and $n=3$ cause
no exception.
\end{proof}
